% ============================================================
%  lang/en.tex --- every user-visible string in the English edition
%
%  Nothing in chapters/en/ may hard-code a word that a macro also typesets.
%  If a label is needed and is not here, add it here AND in lang/pl.tex; CI
%  fails the build when the two files do not define the same macro set (C3).
% ============================================================

% ---------- Part titles ----------
\newcommand{\lblPartI}{The clockwork: how mathematics describes the world}
\newcommand{\lblPartII}{Where prediction breaks}
\newcommand{\lblPartIII}{The shapes of chaos}
\newcommand{\lblPartIV}{Chaos in the world}
\newcommand{\lblPartV}{Living with chaos}
\newcommand{\lblPartAppendices}{Appendices}

% ---------- Admonition titles ----------
\newcommand{\lblNote}{Note}
\newcommand{\lblWarning}{Watch out}
\newcommand{\lblTrap}{The trap}
\newcommand{\lblWorld}{Where you meet this}
\newcommand{\lblRigour}{What we are not proving}
\newcommand{\lblProject}{chaoslab}
\newcommand{\lblVerify}{Run this before you trust it}
\newcommand{\lblLab}{Lab}

% ---------- The method ----------
\newcommand{\lblThink}{Stop and think}
\newcommand{\lblThinkCue}{Answer it before you read on.}
\newcommand{\lblAnswerMark}{Answer.}
\newcommand{\lblOutcomes}{What you will be able to do}
\newcommand{\lblOutcomesLead}{When you have finished this chapter you will be able to:}
\newcommand{\lblSummary}{Summary}
\newcommand{\lblCanYou}{Can you?}
\newcommand{\lblCanYouIntro}{%
  Rate yourself against each item, 1 (no) to 5 (yes, unaided). Anything below
  4 is not a matter for reflection: go back to the section that taught it and
  work its questions again before you move on.}
\newcommand{\lblNoOutcomes}{No outcomes were declared for this chapter.}
\newcommand{\lblCheckpoint}{Checkpoint}
\newcommand{\lblCheckpointIntro}{%
  Answer these on paper before you look at Appendix~A. They test what the
  chapter taught, in the order it taught it.}
\newcommand{\lblAnswersIntro}{%
  The checkpoint answers of every chapter, in chapter order. Each was written
  beside its question and collected here by the build, so an answer cannot
  drift away from the question it answers.}
\newcommand{\lblNoAnswers}{No answers were stored. This is a partial build.}

% ---------- The lab footer ----------
\newcommand{\lblLabStarter}{Starter}
\newcommand{\lblLabTest}{Test}
\newcommand{\lblLabRun}{Run}

% ---------- Build-state markers ----------
\newcommand{\lblNotWritten}{NOT YET WRITTEN}
\newcommand{\lblNotRendered}{DIAGRAM NOT RENDERED --- run \texttt{make diagrams}}
\newcommand{\lblPlotNotRendered}{PLOT NOT RENDERED --- run \texttt{make plots}}
\newcommand{\lblTranscriptMissing}{TRANSCRIPT NOT GENERATED --- run \texttt{make numbers}}
\newcommand{\lblValueMissing}{value missing}
\newcommand{\lblDiagramManifest}{Diagram manifest}
\newcommand{\lblPlotManifest}{Plot manifest}
\newcommand{\lblLabManifest}{Lab manifest}

% ---------- Front-matter running heads ----------
\newcommand{\lblHowToUse}{How to use this book}
\newcommand{\lblIntroduction}{Introduction}

% The punctuation dash, written \dash{} with a space either side and never
% typed directly: the English edition sets a spaced em dash and the Polish
% one a spaced en dash.
\newcommand{\dash}{---}

% ---------- Index cross-references ----------
\newcommand{\lblIndexSee}{see}
\newcommand{\lblIndexSeeAlso}{see also}

% The date the book's facts and pins were last checked.
\newcommand{\pinnedon}{September 2026}

% ---------- PDF metadata ----------
\newcommand{\lblPdfSubject}{How mathematics describes the world, and where it stops predicting it: chaos theory from school arithmetic}
\newcommand{\lblPdfKeywords}{chaos theory, dynamical systems, logistic map, Lorenz attractor, fractals, Lyapunov exponent, prediction, Python}
\newcommand{\lblPdfLang}{en-GB}
